\begin{name}
	{\tenchude}
	{ĐỀ TT LẦN 1 - SỞ HƯNG YÊN}
	{LỚP TOÁN THẦY PHÁT}
	{Um passo à frente e você não está mais no mesmo lugar}
\end{name}
\Opensolutionfile{ans}[ans/ans\currfilebase]
\TN
\begin{ex}%[Đề KSCL lần I, Sở Hưng Yên, năm học 2025-2026]%[Nguyễn Trung Kiên, dự án TeX hóa đề sở năm học 2025-2026]%[2D1N4-1]
	\immini
	{Cho hàm số $y=f(x)$ có đồ thị như hình vẽ. Đồ thị hàm số đã cho có đường tiệm cận đứng là
		\choice[2]
		{$y=2$}
		{$x=2$}
		{\True $x=1$}
		{$x=0$}}
	{\begin{tikzpicture}[scale=1, font=\footnotesize,line join=round, line cap=round, >=stealth, x=0.6cm, y=0.6cm]
			\def\xt{-3} \def\xp{5} \def\yt{5} \def\yd{-3}
			\tikzset{declare function={f(\x)=(3*(\x)+1)/(3*(\x)-3);}}
			\draw[->] (\xt,0)--(\xp,0) node [below]{$x$};
			\draw[->] (0,\yd)--(0,\yt) node [left]{$y$};
			\node at (0,0) [below right]{$O$};
			\begin{scope}
				\clip (\xt,\yd) rectangle (\xp,\yt);
				\draw[smooth,samples=200,domain=\xt:0.9] plot(\x,{f(\x});
				\draw[smooth,samples=200,domain=1.1:\xp] plot(\x,{f(\x});
				\draw[blue] (1,\yd)--(1,\yt);
				\draw[blue] (\xt,1)--(\xp,1);
			\end{scope}			
			\foreach \diem/\pos/\name in {(1,0)/-70/1, (0,1)/130/1} \fill \diem node[shift={(\pos:4mm)}]{$\name$} circle(1pt);
			\fill (-1/3,0) node[shift={(-115:3mm)}]{\scriptsize$-\tfrac{1}{3}$} circle(1pt)
			(0,0) circle(1pt);
	\end{tikzpicture}}
	\loigiai{
		Dựa vào đồ thị, ta thấy đường tiệm cận đứng của đồ thị hàm số đã cho là $x=1$.
	}
\end{ex}

\begin{ex}%[Đề KSCL lần I, Sở Hưng Yên, năm học 2025-2026]%[Nguyễn Trung Kiên, dự án TeX hóa đề sở năm học 2025-2026]%[2H2N1-2]
	Cho hình lập phương $ABCD.A'B'C'D'$ có độ dài cạnh bằng $1$. Độ dài của véc-tơ $\overrightarrow{BA} + \overrightarrow{BC} + \overrightarrow{BB'}$ bằng
	\choice
	{\True $\sqrt{3}$}
	{$2\sqrt{3}$}
	{$3$}
	{$\sqrt{2}$}
	\loigiai{
		Theo quy tắc hình hộp, ta có $\overrightarrow{BA} + \overrightarrow{BC} + \overrightarrow{BB'}=\overrightarrow{BD'}$.\\
		Do đó $\left|\overrightarrow{BA} + \overrightarrow{BC} + \overrightarrow{BB'}\right|=\left|\overrightarrow{BD'}\right|=BD'=\sqrt{3}$.
	}
\end{ex}

\begin{ex}%[Đề KSCL lần I, Sở Hưng Yên, năm học 2025-2026]%[Nguyễn Trung Kiên, dự án TeX hóa đề sở năm học 2025-2026]%[2D3N2-2]
	Xét mẫu số liệu ghép nhóm có phương sai bằng $16$. Độ lệch chuẩn của mẫu số liệu đó bằng
	\choice
	{$2$}
	{\True $4$}
	{$8$}
	{$6$}
	\loigiai{
		Độ lệch chuẩn là căn bậc hai của phương sai $s=\sqrt{s^2}=\sqrt{16}=4$.
	}
\end{ex}

\begin{ex}%[Đề KSCL lần I, Sở Hưng Yên, năm học 2025-2026]%[Nguyễn Trung Kiên, dự án TeX hóa đề sở năm học 2025-2026]%[2H2N2-2]
	Trong không gian với hệ tọa độ $Oxyz$, cho điểm $M$ thỏa mãn $\overrightarrow{OM}=\overrightarrow{i} + 4\overrightarrow{j} - 2\overrightarrow{k}$. Tọa độ của điểm $M$ là
	\choice
	{$(1;4;2)$}
	{\True $(1;4;-2)$}
	{$(-1;4;-2)$}
	{$(-1;-4;2)$}
	\loigiai{
		Vì $\overrightarrow{OM}=\overrightarrow{i} + 4\overrightarrow{j} - 2\overrightarrow{k}$ nên $M(1;4;-2)$.
	}
\end{ex}

\begin{ex}%[Đề KSCL lần I, Sở Hưng Yên, năm học 2025-2026]%[Nguyễn Trung Kiên, dự án TeX hóa đề sở năm học 2025-2026]%[1H8H7-3]
	Cho khối chóp có đáy là hình vuông cạnh bằng $4$ và chiều cao $h=3$. Thể tích $V$ của khối chóp đã cho bằng
	\choice
	{$V=48$}
	{$V=24$}
	{\True $V=16$}
	{$V=8$}
	\loigiai{
		Diện tích đáy $B=4^2=16$ (đvdt).\\
		Thể tích khối chóp $V=\dfrac{1}{3}Bh=\dfrac{1}{3} \cdot 16 \cdot 3=16$ (đvtt).
	}
\end{ex}

\begin{ex}%[Đề KSCL lần I, Sở Hưng Yên, năm học 2025-2026]%[Nguyễn Trung Kiên, dự án TeX hóa đề sở năm học 2025-2026]%[2D1N2-2]
	\immini
	{Cho hàm số $y=f(x)$ có đồ thị như hình vẽ. Điểm cực tiểu của hàm số đã cho là
		\choice[2]
		{\True $x=2$}
		{$x=1$}
		{$x=0$}
		{$x=-2$}}
	{\begin{tikzpicture}[scale=0.9, font=\footnotesize, line join=round, line cap=round, >=stealth]
			\def\xt{-2}
			\def\xp{4}
			\def\yd{-2.5}
			\def\yt{2.5}
			\tikzset{declare function={f(\x)=(\x)^3-3*(\x)^2+2;}}
			\draw[->](\xt,0) -- (0,0)node[below left]{$O$} -- (\xp,0)node[below]{$x$};
			\draw[->](0,\yd)--(0,\yt)node[right]{$y$};
			\begin{scope}
				\clip (\xt,\yd) rectangle (\xp,\yt);
				\draw[smooth,samples=150] plot[domain=\xt:\xp](\x,{f(\x)});
			\end{scope}
			\foreach \i in {(0,0), (2,0), (0,-2), (0,2), (2,-2)} \fill \i circle(1pt);
			\draw[dashed](2,0)node[above]{$2$} |- (0,-2)node[left]{$-2$} 
			(0,2)node[shift={(150:3mm)}]{$2$};
	\end{tikzpicture}}
	\loigiai{
		Dựa vào đồ thị, hàm số đạt cực tiểu tại điểm $x=2$.
	}
\end{ex}

\begin{ex}%[Đề KSCL lần I, Sở Hưng Yên, năm học 2025-2026]%[Nguyễn Trung Kiên, dự án TeX hóa đề sở năm học 2025-2026]%[1D6N4-2]
	Nghiệm của phương trình $\log_4(3x+4)=2$ là
	\choice
	{$x=\dfrac{4}{3}$}
	{\True $x=4$}
	{$x=\dfrac{2}{3}$}
	{$x=8$}
	\loigiai{
		Ta có $$\log_4(3x+4)=2 \Leftrightarrow 3x+4=4^2 \Leftrightarrow 3x+4=16 \Leftrightarrow 3x=12 \Leftrightarrow x=4.$$
	}
\end{ex}

\begin{ex}%[Đề KSCL lần I, Sở Hưng Yên, năm học 2025-2026]%[Nguyễn Trung Kiên, dự án TeX hóa đề sở năm học 2025-2026]%[2H2H2-3]
	Trong không gian với hệ tọa độ $Oxyz$, cho véc-tơ $\overrightarrow{u}=(2; 0; -1)$ và điểm $A(3; 1; -2)$. Tọa độ điểm $B$ thỏa mãn $\overrightarrow{AB}=\overrightarrow{u}$ là
	\choice
	{$(1; 1; 1)$}
	{$(1; -1; 0)$}
	{$(-5; -1; 3)$}
	{\True $(5; 1; -3)$}
	\loigiai{
		Gọi $B(x; y; z)$, ta có $\overrightarrow{AB}=(x-3; y-1; z+2)$.\\
		Do đó $\overrightarrow{AB}=\overrightarrow{u} \Leftrightarrow \heva{&x-3=2 \\ &y-1=0 \\ &z+2=-1} \Leftrightarrow \heva{&x=5 \\ &y=1 \\ &z=-3.}$\\
		Vậy $B(5; 1; -3)$.
	}
\end{ex}

\begin{ex}%[Đề KSCL lần I, Sở Hưng Yên, năm học 2025-2026]%[Nguyễn Trung Kiên, dự án TeX hóa đề sở năm học 2025-2026]%[1D7H2-1]
	Cho hàm số $y=\dfrac{3x-1}{x+3}$ có tập xác định $\mathscr{D}$. Hàm số đã cho có đạo hàm trên $\mathscr{D}$ là
	\choice
	{$y'=-\dfrac{10}{(x+3)^2}$}
	{\True $y'=\dfrac{10}{(x+3)^2}$}
	{$y'=-\dfrac{8}{(x+3)^2}$}
	{$y'=\dfrac{8}{(x+3)^2}$}
	\loigiai
	{Đạo hàm của hàm số $y=\dfrac{3x-1}{x+3}$ trên $(-\infty;-3) \cup (-3;+\infty)$ là 
		$$y'=\dfrac{3(x+3)-(3x-1)}{(x+3)^2}=\dfrac{10}{(x+3)^2}.$$}
\end{ex}

\begin{ex}%[Đề KSCL lần I, Sở Hưng Yên, năm học 2025-2026]%[Nguyễn Trung Kiên, dự án TeX hóa đề sở năm học 2025-2026]%[2D3N1-3]
	Xét mẫu số liệu ghép nhóm có tứ phân vị thứ nhất, tứ phân vị thứ hai và tứ phân vị thứ ba lần lượt là $27{,}5$; $30{,}5$ và $33$. Khoảng tứ phân vị của mẫu số liệu ghép nhóm đó bằng
	\choice
	{$3$}
	{$5$}
	{\True $5{,}5$}
	{$2{,}5$}
	\loigiai
	{Khoảng tứ phân vị của mẫu ghép nhóm là $\Delta_Q=Q_3-Q_1=33-27{,}5=5{,}5$.}
\end{ex}

\begin{ex}%[Đề KSCL lần I, Sở Hưng Yên, năm học 2025-2026]%[Nguyễn Trung Kiên, dự án TeX hóa đề sở năm học 2025-2026]%[2D1H1-1]
	Trong các hàm số cho dưới đây, hàm số nào nghịch biến trên khoảng $(-\infty;+\infty)$?
	\choice
	{$y=\dfrac{1-x^2}{x}$}
	{$y=-x^4-2x^2$}
	{\True $y=-x^3-3x+2$}
	{$y=\dfrac{x+2}{x-1}$}
	\loigiai
	{Ta thấy hàm số $y=\dfrac{1-x^2}{x}$ xác định trên $(-\infty;+\infty)\setminus \{0\}$, hàm số $y=\dfrac{x+2}{x-1}$ xác định trên $(-\infty;+\infty)\setminus \{1\}$ nên không thỏa mãn.\\
		Hàm số $y=-x^3-3x+2$ có tập xác định $(-\infty;+\infty)$ và $y'=-3x^2-3<0,\, \forall x\in (-\infty;+\infty)$ nên nghịch biến trên khoảng $(-\infty;+\infty)$.}
\end{ex}

\begin{ex}%[Đề KSCL lần I, Sở Hưng Yên, năm học 2025-2026]%[Nguyễn Trung Kiên, dự án TeX hóa đề sở năm học 2025-2026]%[2D3N1-2]
	Kết quả đo chiều cao của $45$ học sinh lớp 12A được thống kê như sau
	\begin{center}
		{\setlength\extrarowheight{2pt}
			\begin{tabular}{|>{\centering\arraybackslash}m{3.5cm}|*{6}{c|}}
				\hline 
				Chiều cao (cm)	& $[150;155)$ & $[155;160)$ & $[160;165)$ & $[165;170)$ & $[170;175)$ & $[175;180)$\\
				\hline 
				Số học sinh	& $6$ & $12$ & $13$ & $8$ & $4$ & $2$\\
				\hline 
		\end{tabular}}
	\end{center}
	Khoảng biến thiên của mẫu số liệu ghép nhóm trên bằng
	\choice
	{\True $30$}
	{$35$}
	{$20$}
	{$25$}
	\loigiai
	{Khoảng biến thiên của mẫu số liệu ghép nhóm trên là $R=180-150=30$.}
\end{ex}
\cauds
\begin{ex}%[Đề thi KSCL học kỳ I, Sở Hưng Yên năm học 2025-2026]%[2D1V5-2]
	Cho hàm số $y=f(x)=\dfrac{2x-1}{x-1}$ có đồ thị là đường cong $(C)$.
	\choiceTF
	{\True Đồ thị $(C)$ như hình vẽ dưới đây 
		\begin{center}  
			\begin{tikzpicture}[scale=1, font=\footnotesize,line join=round, line cap=round, >=stealth, x=0.5cm, y=0.5cm]
				\def\xt{-3.5} \def\xp{5.5} \def\yt{6} \def\yd{-2}
				\tikzset{declare function = {f(\x)= (2*(\x)-1)/(\x-1);}}
				\draw[->] (\xt,0)--(\xp,0) node [below]{$x$};
				\draw[->] (0,\yd)--(0,\yt) node [left]{$y$};
				\node at (0,0) [below left]{$O$};
				\begin{scope}
					\clip (\xt,\yd) rectangle (\xp,\yt);
					\draw[smooth,samples=200,domain=\xt:0.9] plot(\x,{f(\x});
					\draw[smooth,samples=200,domain=1.1:\xp] plot(\x,{f(\x});
					\draw[blue] (1,\yd)--(1,\yt);
					\draw[blue] (\xt,2)--(\xp,2);
				\end{scope}			
				\foreach \diem/\pos/\name in {(1,0)/-65/1, (2,0)/-90/2, (0,1)/190/1, (0,2)/145/2, (0,3)/170/3} \fill \diem node[shift={(\pos:3mm)}]{$\name$} circle(1pt);
				\draw[dashed] (2,0)|-(0,3);
				\fill (0,0) circle(1pt) (2,3) circle(1pt);
			\end{tikzpicture}
	\end{center}}
	{Hàm số nghịch biến trên khoảng $(-\infty; +\infty)$}
	{\True Đồ thị $(C)$ có đường tiệm cận ngang $y=2$ và đường tiệm cận đứng $x=1$}
	{Gọi $M$, $m$ lần lượt là giá trị lớn nhất, giá trị nhỏ nhất của hàm số $y=|f(x)|$ trên đoạn $\left[-1; \dfrac{2}{3}\right]$. Khi đó $2M + 2\,023m=2\,026$}
	\loigiai{
		\begin{itemchoice}
			\itemch Đồ thị hàm số có đường tiệm cận đứng $x=1$, đường tiệm cận ngang $y=2$, đi qua các điểm $(0;1)$ và $(2;3)$.
			\itemch Tập xác định $\mathscr{D}=\mathbb{R}\setminus \{1\}$.\\
			Ta có $y'=\dfrac{-1}{(x-1)^2}<0$, $\forall x\in\mathscr{D}$.\\
			Vậy hàm số đã cho nghịch biến trên các khoảng $(-\infty;1)$ và $(1;+\infty)$.
			\itemch Đồ thị $(C)$ có đường tiệm cận ngang $y=2$ và đường tiệm cận đứng $x=1$.
			\itemch Xét đồ thị của hàm số $y=|f(x)|$ trên đoạn $\left[-1; \dfrac{2}{3}\right]$ như sau
			\begin{center}
				\begin{tikzpicture}[scale=1, font=\footnotesize,line join=round, line cap=round, >=stealth, x=1.2cm, y=1.2cm]
					\def\xt{-2} \def\xp{2} \def\yt{3} \def\yd{-0.7}
					\tikzset{declare function = {f(\x)= (2*(\x)-1)/(\x-1);}}
					\draw[->] (\xt,0)--(\xp,0) node [below]{$x$};
					\draw[->] (0,\yd)--(0,\yt) node [left]{$y$};
					\node at (0,0) [below left]{$O$};
					\begin{scope}
						\clip (\xt,\yd) rectangle (\xp,\yt);
						\draw[smooth,samples=200,domain=\xt:0.5] plot(\x,{f(\x});
						\draw[smooth,samples=200,domain=0.5:0.9] plot(\x,{-f(\x});
						\draw[blue] (1,\yd)--(1,\yt);
					\end{scope}		
					\draw[dashed] (-1,0) node[below]{$-1$} --(-1,1.5)--(0,1.5) node[above left]{$\dfrac{3}{2}$};
					\draw[dashed] (2/3,0) node[shift={(-80:5mm)}]{$\dfrac{2}{3}$}|-(0,1) node[shift={(190:3mm)}]{$1$};
					\draw (0.5,0) node[shift={(-100:5mm)}]{$\dfrac{1}{2}$};
					\foreach \i in {(-1,3/2), (0,0), (-1,0), (1/2,0), (2/3,0), (0,1), (2/3,1), (0,3/2)} \fill \i circle(1pt); 
				\end{tikzpicture}	
			\end{center}
			Dựa vào đồ thị ta thấy $M=\dfrac{3}{2}$, $m=0$.\\
			Vậy $2M+2\,023 m=3$.
		\end{itemchoice}	
	}
\end{ex}

\begin{ex}%[Đề thi KSCL học kỳ I, Sở Hưng Yên năm học 2025-2026]%[1D9H2-4]
	Phỏng vấn $30$ học sinh lớp 11A của một trường THPT về môn thể thao yêu thích, thu được kết quả: Có $18$ học sinh thích bóng đá, $15$ học sinh thích bóng bàn và $9$ học sinh thích cả hai môn thể thao bóng đá và bóng bàn. Chọn ngẫu nhiên một học sinh trong số các học sinh đã phỏng vấn.
	\choiceTF
	{Xác suất để chọn được học sinh chỉ thích một môn bóng bàn bằng $0{,}5$}
	{\True Xác suất để chọn được học sinh thích cả hai môn bóng đá và bóng bàn bằng $0{,}3$}
	{\True Xác suất để chọn được học sinh không thích môn thể thao nào trong hai môn trên bằng $0{,}2$}
	{\True Xác suất để chọn được học sinh thích bóng đá bằng $0{,}6$}
	\loigiai{
		Gọi $X$ là tập hợp học sinh lớp 11A tham gia phỏng vấn; $A$ là tập hợp học sinh thích bóng đá; $B$ là tập hợp học sinh thích bóng bàn.
		\begin{center}
			\begin{tikzpicture}[scale=1, font=\footnotesize, line join=round, line cap=round, >=stealth]
				\draw (0,0) ellipse (2.7 and 1.6);
				\draw[name path=A, fill=cyan!30] (-0.5,-0.15) ellipse (1.5 and 1);
				\draw[name path=B, pattern=dots, pattern color=red] (0.75,0.1) ellipse (1.5 and 1);
				\draw[<-] (180:2.7 and 1.6) --++(180:0.5) node[left]{$X$};
				\draw[<-] ($(-0.5,-0.15)+(220:1.5 and 1)$) --++(200:1.5) node[left]{$A$};
				\draw[<-] ($(0.75,0.1)+(-30:1.5 and 1)$) --++(-20:1.5) node[right]{$B$};
				\draw[<-] (160: 2.3 and 1.4) --++(160:1.5) node[left]{$\overline{A} \cap \overline{B}$};
				\path[name intersections = {of=A and B, by={M,N}}] ($(M)!1/2!(N)$) node{$A\cap B$};
			\end{tikzpicture}
		\end{center}
		Ta có $n(X)=30$, $n(A)=18$, $n(B)=15$, $n(AB)=9$.
		\begin{itemchoice}
			\itemch Số học sinh chỉ thích một môn bóng bàn là $15-9=6$.\\
			Vậy xác suất để chọn được học sinh chỉ thích một môn bóng bàn bằng $\dfrac{6}{30}=0{,}2$.
			\itemch Xác suất để chọn được học sinh thích cả hai môn bóng đá và bóng bàn bằng $\dfrac{9}{30}=0{,}3$.
			\itemch Số học sinh thích ít nhất một trong hai môn bóng đá hoặc bóng bàn là $$n(A\cup B) = n(A)+n(B)-n(AB) = 18+15-9=24 \text{ (học sinh)}.$$
			Số học sinh không thích môn thể thao nào trong hai môn trên là $30-24=6$ (học sinh).\\
			Vậy xác suất để chọn được học sinh không thích bóng đá và cũng không thích bóng bàn bằng $\dfrac{6}{30}=0{,}2$.
			\itemch Xác suất để chọn được học sinh thích bóng đá bằng $\dfrac{18}{30}=0{,}6$.
		\end{itemchoice}	
	}
\end{ex}

\begin{ex}%[Đề thi KSCL học kỳ I, Sở Hưng Yên năm học 2025-2026]%[2H2V2-4]
	\immini
	{Cho hình chóp tứ giác đều $S.ABCD$, $O$ là tâm của đáy $ABCD$ được gắn vào hệ trục tọa độ $Oxyz$ như hình vẽ. Biết cạnh $SA=AB=3\sqrt{2}$ và điểm $G$ là trọng tâm tam giác $SAB$. Khi đó
		\choiceTF
		{Tọa độ điểm $B$ là $(6;0;0)$}
		{Độ dài đoạn $DG$ bằng $18$}
		{\True $\overrightarrow{GA}+\overrightarrow{GS}+\overrightarrow{GB}=\overrightarrow{0}$}
		{\True Nếu $K(0;m;n)$ là điểm thuộc mặt phẳng $(Oyz)$ sao cho ${KG+KB}$ đạt giá trị nhỏ nhất thì $3m^2+4n^2=\dfrac{63}{16}$}	
	}
	{\begin{tikzpicture}[scale=1, font=\footnotesize,line join=round, line cap=round, >=stealth]
			\def\h{3}
			\path (0,0) coordinate (A)
			(-1.6,-1.2) coordinate (B)
			++(4,0) coordinate (C)
			($(A)+(C)-(B)$) coordinate(D)
			($(A)!1/2!(C)$) coordinate(O)
			++(0,\h) coordinate(S)
			($1/3*(A)+1/3*(B)+1/3*(S)$) coordinate(G)
			;
			\draw (S)--(B)--(C)--(D)--(S)--(C);
			\draw[dashed] (S)--(A)--(B)--(D)--(A)--(C) (S)--(O);
			\foreach \diem/\pos in {A/170, B/-90, C/-120, D/0, S/160,O/-90,G/60} \fill (\diem) node[shift={(\pos:0.3)}] {$\diem$} circle(1pt);
			\draw[->] (B)--($(O)!1.3!(B)$) node[above]{$x$};
			\draw[->] (C)--($(O)!1.3!(C)$) node[right]{$y$};
			\draw[->] (S)--($(O)!1.3!(S)$) node[left]{$z$};
	\end{tikzpicture}}
	\loigiai{
		Do $S.ABCD$ là hình chóp đều nên đáy $ABCD$ là hình vuông tâm $O$ và $SO\perp (ABCD)$.\\
		Do $AB=3\sqrt{2}$ nên $OA=OB=OC=OD=3$.\\
		Tam giác $SOA$ vuông tại $O$ nên $SO=\sqrt{SA^2-AO^2}=3$.\\
		Với việc gắn hệ tọa độ như hình đã cho, ta có $A(0;-3;0)$, $B(3;0;0)$, $S(0;0;3)$, $D(-3;0;0)$.
		\begin{itemchoice}
			\itemch Tọa độ điểm $B$ là $(3;0;0)$.
			\itemch Với $G$ là trọng tâm của tam giác $SAB$ thì $\heva{& x_G=\dfrac{x_A+x_B+x_S}{3}=1 \\ & y_G=\dfrac{y_A+y_B+y_S}{3}=-1\\ &z_G=\dfrac{z_A+z_B+z_S}{3}=1.}$\\
			Vậy $G(1;-1;1)$, với $D(-3;0;0)$ thì $DG=\sqrt{4^2+(-1)^2+1^2}=3\sqrt{2}$.
			\itemch Với $G$ là trọng tâm tam giác $SAB$ thì $\overrightarrow{GA}+\overrightarrow{GS}+\overrightarrow{GB}=\overrightarrow{0}$
			\itemch Xét hình vẽ sau
			\begin{center}
				\begin{tikzpicture}[scale=1, font=\footnotesize,line join=round, line cap=round, >=stealth]
					\def\h{3}
					\path (0,0) coordinate (A)
					(-1.6,-1.2) coordinate (B)
					++(4,0) coordinate (C)
					($(A)+(C)-(B)$) coordinate(D)
					($(A)!1/2!(C)$) coordinate(O)
					++(0,\h) coordinate(S)
					($1/3*(A)+1/3*(B)+1/3*(S)$) coordinate(G)
					($(S)!0.6!(O)$) coordinate(K)
					;
					\draw (S)--(B)--(C)--(D)--(S)--(C);
					\draw[dashed] (S)--(A)--(B)--(D)--(A)--(C) (S)--(O) (B)--(K)--(D) (K)--(G) (G)--(D);
					\foreach \diem/\pos in {A/170, B/-90, C/-120, D/0, S/160,O/-90,G/60,K/120} \fill (\diem) node[shift={(\pos:0.3)}] {$\diem$} circle(1pt);
					\draw[->] (B)--($(O)!1.3!(B)$) node[above]{$x$};
					\draw[->] (C)--($(O)!1.3!(C)$) node[right]{$y$};
					\draw[->] (S)--($(O)!1.3!(S)$) node[left]{$z$};
				\end{tikzpicture}	
			\end{center}
			Ta có $D$ là điểm đối xứng với $B$ qua $(Oyz)$.\\ 
			Khi đó $KB+KG=KD+KG\ge DG=3\sqrt{2}$.\\
			Dấu bằng xảy ra khi ba điểm $D$, $K$, $G$ thẳng hàng và $K$ nằm giữa $D$ và $G$.\\
			Ta có $\overrightarrow{DG}=(4;-1;1)$, $\overrightarrow{DK}=(3;m;n)$.\\
			Hai véc-tơ $\overrightarrow{DG}$ và $\overrightarrow{DK}$ cùng phương suy ra 
			$$\dfrac{3}{4}=\dfrac{m}{-1}=\dfrac{n}{1}\Leftrightarrow \heva{& m=-\dfrac{3}{4} \\ & n=\dfrac{3}{4}.}$$
			Vậy $3m^2+4n^2=\dfrac{63}{16}$.
		\end{itemchoice}	
	}
\end{ex}

\begin{ex}%[Đề thi KSCL học kỳ I, Sở Hưng Yên năm học 2025-2026]%[2D3H2-2]
	Kết quả điểm thi khảo sát chất lượng học kỳ I môn Toán hai khối $10$; $11$ ở một trường THPT tỉnh Hưng Yên được biểu diễn bởi mẫu số liệu ghép nhóm ở bảng sau:
	\begin{center}
		{\setlength\extrarowheight{2pt}
			\begin{tabular}{|>{\centering\arraybackslash}m{2.5cm}|*{9}{c|}}
				\hline 
				{\bf Nhóm} & $[1;2)$ & $[2;3)$ & $[3;4)$ & $[4;5)$ & $[5;6)$ & $[6;7)$ & $[7;8)$ & $[8;9)$ & $[9;10]$\\
				\hline 
				{\bf Khối 11}& $11$ & $52$ & $118$ & $121$ & $95$ & $62$ & $47$ & $27$ & $7$\\
				\hline 
				{\bf Khối 10} & $1$ & $4$ & $35$ & $86$ & $117$ & $112$ & $102$ & $64$ & $19$\\
				\hline 
		\end{tabular}}
	\end{center}
	\choiceTF
	{\True Khoảng biến thiên của mẫu số liệu ghép nhóm là $R=9$}
	{Điểm trung bình của toàn khối $10$ là $6{,}13$ (kết quả làm tròn đến hàng phần trăm)}
	{\True Nếu so sánh theo phương sai của mẫu số liệu ghép nhóm thì học sinh khối $10$ có điểm đồng đều hơn điểm của học sinh khối $11$}
	{\True Khoảng tứ phân vị của mẫu số liệu ghép nhóm về điểm thi của khối $11$ là $\Delta_{Q_{K11}}=2{,}52$ (kết quả tính làm tròn đến hàng phần trăm)}
	\loigiai
	{Ta có bảng số liệu ghép nhóm bổ sung giá trị đại diện như sau:
		\begin{center}
			{\setlength\extrarowheight{2pt}
				\begin{tabular}{|>{\centering\arraybackslash}m{3cm}|*{9}{c|}}
					\hline 
					{\bf Nhóm} & $[1;2)$ & $[2;3)$ & $[3;4)$ & $[4;5)$ & $[5;6)$ & $[6;7)$ & $[7;8)$ & $[8;9)$ & $[9;10]$\\
					\hline 
					{\bf Giá trị đại diện} &$1{,}5$ & $2{,}5$ & $3{,}5$ & $4{,}5$ & $5{,}5$ & $6{,}5$ & $7{,}5$ & $8{,}5$ & $9{,}5$\\
					\hline
					{\bf Khối 11}& $11$ & $52$ & $118$ & $121$ & $95$ & $62$ & $47$ & $27$ & $7$\\
					\hline 
					{\bf Khối 10} & $1$ & $4$ & $35$ & $86$ & $117$ & $112$ & $102$ & $64$ & $19$\\
					\hline 
			\end{tabular}}	
		\end{center}
	\noindent
		Số học sinh của khối $11$ là $n=11+52+118+121+95+62+47+27+7=540$.\\
		Số học sinh của khối $10$ là $n=1+4+35+86+117+112+102+64+19=540$.
		\begin{itemchoice}
			\itemch Khoảng biến thiên của mẫu số liệu ghép nhóm là $R=10-1=9$.
			\itemch Điểm trung bình của toàn khối $10$ là
			$$\overline{x}_{10}=\dfrac{1{,}5\cdot 1+2{,}5\cdot 4+3{,}5\cdot 35+\cdots +9{,}5\cdot 19}{540}=\dfrac{1\,691}{270}\approx 6{,}26.$$
			\itemch 
			Phương sai của mẫu số liệu ghép nhóm khối $10$ là
			$$s^2_{10}=\dfrac{1{,}5^2\cdot 1+2{,}5^2\cdot 4+3{,}5^2\cdot 35+\cdots +9{,}5^2\cdot 19}{540}-\left(\overline{x}_{10}\right)^2=\dfrac{22\,545}{540}-\left(\dfrac{1\,691}{270}\right)^2\approx 2{,}53.$$
			Điểm trung bình của toàn khối $11$ là
			$$\overline{x}_{11}=\dfrac{1{,}5\cdot 11+2{,}5\cdot 52+3{,}5\cdot 118+\cdots +9{,}5\cdot 7}{540}=\dfrac{1\,339}{270}.$$
			Phương sai của mẫu số liệu ghép nhóm khối $11$ là
			$$s^2_{11}=\dfrac{1{,}5^2\cdot 11+2{,}5^2\cdot 52+3{,}5^2\cdot 118+\cdots +9{,}5^2\cdot 7}{540}-\left(\overline{x}_{11}\right)^2=\dfrac{14\,965}{540}-\left(\dfrac{1\,339}{270}\right)^2\approx 3{,}12.$$
			Vậy nếu so sánh theo phương sai của mẫu số liệu ghép nhóm thì học sinh khối $10$ có điểm đồng đều hơn điểm của học sinh khối $11$	
			\itemch 
			Ta có $\dfrac{n}{4}=135$, $11+52+118>135$ nên nhóm chứa $Q_1$ là $[3;4)$.\\
			Ta có $Q_1=3+\dfrac{\dfrac{n}{4}-(11+52)}{118}\cdot (4-3)=\dfrac{213}{59}$.\\
			Ta có $\dfrac{3n}{4}=405$, và $11+52+118+121+95+62>405$ nên nhóm chứa $Q_3$ là $[6;7)$.\\
			Ta có $Q_3=6+\dfrac{\dfrac{3n}{4}-(11+52+118+121+95)}{62}\cdot (7-6)=\dfrac{190}{31}$.\\
			Vậy khoảng tứ phân vị của mẫu số liệu ghép nhóm về điểm thi của khối $11$ là $$\Delta_{Q_{K11}}= Q_3-Q_1=\dfrac{190}{31}-\dfrac{213}{59} = \dfrac{4\,607}{1\,829} \approx 2{,}52$$
			
			
		\end{itemchoice}	
	}
\end{ex}
\caukq
\begin{ex}%[Đề KSCL lần I, Sở Hưng Yên, năm học 2025-2026]%[2D1V3-6]
	Doanh số (tính bằng số sản phẩm) của một sản phẩm mới (trong vòng một số năm nhất định) tuân theo quy luật logistic được mô hình hóa bằng hàm số $f(t)=\dfrac{24\,000}{1+6\mathrm{e}^{-t}}$ với $t \geq 0$, trong đó thời gian $t$ được tính bằng năm, kể từ khi phát hành sản phẩm mới. Khi đó, đạo hàm $f'(t)$ sẽ biểu thị tốc độ bán hàng. Tốc độ bán hàng lớn nhất đạt được khi $t=\ln a$. Tìm $a$.
	\shortans{6}
	\loigiai{
		Tốc độ bán hàng
		\begin{eqnarray*}
			f'(t) &=& 24\,000 \cdot \left[ \left(1 + 6\mathrm{e}^{-t}\right)^{-1} \right]' \\
			&=& 24\,000 \cdot (-1)\left(1 + 6\mathrm{e}^{-t}\right)^{-2} \cdot \left(-6\mathrm{e}^{-t}\right) \\
			&=& \dfrac{144\,000\mathrm{e}^{-t}}{\left(1 + 6\mathrm{e}^{-t}\right)^2}.
		\end{eqnarray*}
		Ta có 
		\begin{eqnarray*}
			f''(t) &=& 144\,000 \cdot \dfrac{-\mathrm{e}^{-t}\left(1 + 6\mathrm{e}^{-t}\right)^2 - \mathrm{e}^{-t} \cdot 2\left(1 + 6\mathrm{e}^{-t}\right) \cdot \left(-6\mathrm{e}^{-t}\right)}{\left(1 + 6\mathrm{e}^{-t}\right)^4} \\
			&=& 144\,000 \cdot \dfrac{-\mathrm{e}^{-t}\left(1 + 6\mathrm{e}^{-t}\right) + 12\mathrm{e}^{-2t}}{\left(1 + 6\mathrm{e}^{-t}\right)^3} \\
			&=& 144\,000 \cdot \dfrac{6\mathrm{e}^{-2t} - \mathrm{e}^{-t}}{\left(1 + 6\mathrm{e}^{-t}\right)^3}.
		\end{eqnarray*}
		Xét phương trình $f''(t) = 0$, ta có 
		$$6\mathrm{e}^{-2t} - \mathrm{e}^{-t} = 0 \Leftrightarrow \mathrm{e}^{-t}(6\mathrm{e}^{-t} - 1) = 0.$$
		Vì $\mathrm{e}^{-t} > 0$ với mọi $t$, nên ta có
		$$6\mathrm{e}^{-t} - 1 = 0 \Leftrightarrow \mathrm{e}^{-t} = \dfrac{1}{6} \Leftrightarrow -t = \ln\left(\dfrac{1}{6}\right) \Leftrightarrow t = \ln 6.$$
		Bảng biến thiên của hàm số $f'(t)$ như sau
		\begin{center}
			\begin{tikzpicture}
				\tkzTabInit[nocadre=false,lgt=1.2,espcl=3.5,deltacl=0.6]
				{$t$ /0.6,$f''(t)$ /0.6,$f'(t)$ /2}
				{$0$,$2$,$+\infty$}
				\tkzTabLine{,+,$0$,-,}
				\tkzTabVar{-/$f'(0)$,+/$6\,000$,-/$0$}
			\end{tikzpicture}
		\end{center}
		Vậy tốc độ bán hàng lớn nhất đạt tại $t=\ln{6}$, suy ra $a = 6$.
	}
\end{ex}

\begin{ex}%[Đề KSCL lần I, Sở Hưng Yên, năm học 2025-2026]%[2H2V2-4]
	Cho hình chóp $S.ABC$ có đáy là tam giác đều có cạnh bằng $2$ và $SA \perp (ABC)$. Gọi $M$, $N$ là hai điểm lần lượt thuộc các cạnh $SB$, $SC$ sao cho $SM=3MB$, $NC=2NS$. Biết rằng $AN$ vuông góc $CM$. Tính thể tích khối chóp $S.ABC$ (làm tròn kết quả đến hàng phần trăm).
	\shortans{1,29}
	\loigiai{
		\begin{center}
			\begin{tikzpicture}[scale=1.1, font=\footnotesize, line join=round, line cap=round, >=stealth]
				\path
				(0,0) coordinate (A)
				(4,0) coordinate (C)
				(2,-2) coordinate (B)
				(0,4) coordinate (S)
				;
				%Lấy H là trung điểm của AB
				\coordinate (H) at ($(A)!0.5!(B)$);
				%Lấy x, y, z
				\coordinate (X) at ($(A)!1.2!(B)$);
				\coordinate (Y) at ($(A)!1.2!(C)$);
				\coordinate (vAS) at ($(S)-(A)$);
				\coordinate (Z)   at ($(H) + 1.2*(vAS)$);
				%Vẽ các tia tọa độ
				\draw[->] (B)--(X) node[right] {$x$};
				\draw[->] (C)--(Y) node[above] {$y$};
				\draw[->] (H)--(Z) node[above] {$z$};
				%Lấy M, N
				\coordinate (M) at ($(S)!3/4!(B)$);    % SM = 3MB
				\coordinate (N) at ($(S)!1/3!(C)$);    % NC = 2NS 
				%Đánh nhãn điểm
				\foreach \x/\g in {A/180,B/-135,C/45,H/-135,S/90,M/180,N/20}\fill[black](\x) circle (1pt) +(\g:3mm) node {$\x$};
				\draw (S)--(A) (S)--(B) (S)--(C) (A)--(B) (B)--(C) (C)--(M);
				\draw[dashed] (C)--(H) (A)--(C) (A)--(N); 
				% Đánh dấu góc vuông 
				\foreach \x/\y/\z in {Z/H/C, B/H/C, S/A/B, S/A/C}{\draw pic[draw,angle radius=1.5mm]{right angle=\x--\y--\z};}
			\end{tikzpicture}
		\end{center}
		Gọi $H$ là trung điểm của $AB$, ta có $HC\perp AB$ và $HC=\sqrt{3}$.\\ 
		Đặt $h$ ($h>0$) là độ dài của đoạn $AS$.\\
		Chọn hệ trục tọa độ $Hxyz$ với gốc tọa độ tại $H(0;0;0)$; tia $Hx$ trùng tia $HB$; tia $Hy$ trùng tia $HC$; tia $Hz$ song song với $AS$.\\
		Tọa độ các đỉnh: $A(-1; 0; 0)$, $B(1; 0; 0)$, $C\left(0; \sqrt{3}; 0\right)$, $S(-1; 0; h)$.\\
		Điểm $M$ thuộc đoạn $SB$ thỏa mãn $SM = 3MB \Leftrightarrow \overrightarrow{SM} = 3\overrightarrow{MB}$, do đó
			\begin{eqnarray*}
				\heva{& x_M+1=3(1-x_M)\\& y_M=3(-y_M)\\& z_M-h=3(-z_M)} \Leftrightarrow \heva{& x_M=\dfrac{1}{2}\\& y_M=0\\& z_M=\dfrac{h}{4}} \Rightarrow M\left(\dfrac{1}{2}; 0; \dfrac{h}{4}\right).
			\end{eqnarray*}
			Điểm $N$ thuộc đoạn $SC$ thỏa mãn $NC = 2NS \Leftrightarrow \overrightarrow{CN} = 2\overrightarrow{NS}$, do đó
			\begin{eqnarray*}
				\heva{& x_N=2(-1-x_N)\\& y_N-\sqrt{3}=2(-y_N)\\& z_N=2(h-z_N)} \Leftrightarrow \heva{& x_N=-\dfrac{2}{3}\\& y_N=\dfrac{\sqrt{3}}{3}\\& z_N=\dfrac{2h}{3}} \Rightarrow N\left(-\dfrac{2}{3}; \dfrac{\sqrt{3}}{3}; \dfrac{2h}{3}\right).
			\end{eqnarray*}
			Từ đó ta có tọa độ các véc-tơ $\overrightarrow{AN} = \left(\dfrac{1}{3}; \dfrac{\sqrt{3}}{3}; \dfrac{2h}{3}\right)$, $\overrightarrow{CM} = \left(\dfrac{1}{2}; -\sqrt{3}; \dfrac{h}{4}\right)$.\\
		Vì $AN \perp CM$ nên $\overrightarrow{AN} \cdot \overrightarrow{CM} = 0$, do đó
		\[ \dfrac{1}{3} \cdot \dfrac{1}{2} + \dfrac{\sqrt{3}}{3} \cdot \left(-\sqrt{3}\right) + \dfrac{2h}{3} \cdot \dfrac{h}{4} = 0 \Leftrightarrow h = \sqrt{5}.\]
		Tam giác $ABC$ là tam giác đều có cạnh bằng $2$ nên có diện tích là
		\[ S_{ABC} = \dfrac{2^2 \cdot \sqrt{3}}{4} = \sqrt{3}.\]
		Chiều cao của khối chóp là $h = SA = \sqrt{5}$.
		Thể tích khối chóp $S.ABC$ là:
		\[ V = \dfrac{1}{3} \cdot S_{ABC} \cdot h = \dfrac{1}{3} \cdot \sqrt{3} \cdot \sqrt{5} = \dfrac{\sqrt{15}}{3} \approx 1{,}29.\]
	}
\end{ex}

\begin{ex}%[Đề KSCL lần I, Sở Hưng Yên, năm học 2025-2026]%[2H2H2-3]
	Trong không gian với hệ tọa độ $Oxyz$, cho ba điểm $A$, $B$, $C$ biết $\overrightarrow{AB}=(3; -3; 3)$, $\overrightarrow{AC}=(-3; 3; -9)$. Điểm $M$ thuộc đoạn $BC$ thỏa mãn $BM=2MC$. Gọi $(a; b; c)$ là tọa độ của $\overrightarrow{AM}$. Tính $b - 2\,030a + c$.
	\shortans{2026}
	\loigiai{
		Điểm $M$ thuộc đoạn $BC$ thỏa mãn $BM=2MC$ nên $\overrightarrow{BM} = 2\overrightarrow{MC}$.
		Ta có
		\begin{eqnarray*}
			\overrightarrow{AM} - \overrightarrow{AB} = 2\left(\overrightarrow{AC} - \overrightarrow{AM}\right) &\Leftrightarrow& 3\overrightarrow{AM} = \overrightarrow{AB} + 2\overrightarrow{AC} \\
			&\Leftrightarrow& \overrightarrow{AM} = \dfrac{1}{3}\overrightarrow{AB} + \dfrac{2}{3}\overrightarrow{AC}.
		\end{eqnarray*}		
		Tọa độ của $\overrightarrow{AM} = (a; b; c)$ được tính như sau
		\[
		\heva{& a = \dfrac{1}{3} \cdot 3 + \dfrac{2}{3} \cdot (-3) = 1 - 2 = -1, \\
			& b = \dfrac{1}{3} \cdot (-3) + \dfrac{2}{3} \cdot 3 = -1 + 2 = 1, \\
			& c = \dfrac{1}{3} \cdot 3 + \dfrac{2}{3} \cdot (-9) = 1 - 6 = -5.}
		\]		
		Giá trị của biểu thức là
		\begin{eqnarray*}
			b - 2\,030a + c &=& 1 - 2\,030 \cdot (-1) - 5 \\
			&=& 1 + 2\,030 - 5 \\
			&=& 2\,026.
		\end{eqnarray*}
	}
\end{ex}

\begin{ex}%[Đề KSCL lần I, Sở Hưng Yên, năm học 2025-2026]%[2D1V2-1]
	Cho hàm số $y=\dfrac{x^2+x+1}{x+1}$ có đồ thị $(C)$ và điểm $M(1; 1)$. Gọi $A$, $B$ là hai điểm cực trị của đồ thị $(C)$. Tính diện tích của tam giác $MAB$.
	\shortans{2}
	\loigiai{
		Tập xác định $\mathscr{D} = \mathbb{R}\setminus \{-1\}$.\\
		Hàm số đã cho có đạo hàm $y' = \dfrac{x^2 + 2x}{(x+1)^2},\, (x\neq -1)$.\\
		Xét phương trình $y' = 0$
		\[ x^2 + 2x = 0 \Leftrightarrow \hoac{& x = 0 \Rightarrow y = 1 \\ & x = -2 \Rightarrow y = -3.}\]
		Tọa độ hai điểm cực trị là $A(0; 1)$ và $B(-2; -3)$.
		Với điểm $M(1; 1)$, ta có $\overrightarrow{MA}=(-1;0)$, $\overrightarrow{MB}=(-3;-4)$.\\
		Ta có $\cos\widehat{AMB} = \dfrac{\overrightarrow{MA} \cdot \overrightarrow{MB}}{MA\cdot MB} = \dfrac{3}{1\cdot 5} = \dfrac{3}{5}$.\\
		Do $0^\circ < \widehat{AMB} <180^\circ$ nên $\sin\widehat{AMB}>0$, $\sin\widehat{AMB} = \sqrt{1-\left(\cos\widehat{AMB}\right)^2} = \dfrac{4}{5}$.\\
		Vậy diện tích tam giác $MAB$ là $S_{MAB}=\dfrac{1}{2}MA\cdot MB\cdot \sin\widehat{AMB} = \dfrac{1}{2}\cdot 1\cdot 5 \cdot \dfrac{4}{5} = 2$.\\
		\textit{Cách khác}. Tính nhanh diện tích tam giác $MAB$ với $\overrightarrow{MA}=(-1;0)$, $\overrightarrow{MB}=(-3;-4)$ như sau
		\[ S_{\triangle MAB} = \dfrac{1}{2} \cdot \left| (-1) \cdot (-4) - (-3) \cdot 0 \right| = 2.\]
	}
\end{ex}

\begin{ex}%[Đề KSCL lần I, Sở Hưng Yên, năm học 2025-2026]%[1C2H3-1]
	Nam và ba người bạn lên kế hoạch cho một chuyến đi phượt xuyên Việt. Lịch trình xuất phát từ Hà Nội (HN), đi qua các thành phố Đà Nẵng (ĐN), Thành phố Hồ Chí Minh (HCM), Cần Thơ (CT) đúng một lần rồi trở về Hà Nội. Sơ đồ chi phí nhiên liệu (tính bằng lít xăng) để đi qua lại giữa các thành phố được biểu diễn như hình vẽ dưới đây. 
	\begin{center}
		\begin{tikzpicture}[scale=1, font=\footnotesize, line join=round, line cap=round, >=stealth]
			\draw (0,0) coordinate(A) -- (2.5,-1.8) coordinate(B) node[midway, below, sloped]{$40$} -- (6,0) coordinate(C) node[midway, below, sloped]{$70$} -- (4.5,1.5) coordinate(D) node[midway, above, sloped]{$30$} -- (A) node[midway, above, sloped]{$90$} -- (C) node[rounded corners, midway, left, fill=cyan, text=black, sloped]{$110$} (B)--(D) node[rounded corners, midway, left, fill=cyan, text=black, sloped]{$50$};
			\foreach \i/\j/\k in {A/HN/180, B/ĐN/-90, C/CT/0, D/HCM/90} \fill(\i) node[shift={(\k:5mm)}]{\j} circle(2pt);
		\end{tikzpicture}
	\end{center}	
	Để hoàn thành hành trình, Nam và nhóm bạn cần đổ ít nhất bao nhiêu lít xăng?
	\shortans{230}
	\loigiai{
		Các hành trình đi qua mỗi thành phố đúng một lần và trở về Hà Nội gồm có
		\begin{itemize}
			\item Hành trình HN $\to$ ĐN $\to$ CT $\to$ HCM $\to$ HN có chi phí
			\[ L_1 = 40 + 70 + 30 + 90 = 230.\]
			\item Hành trình HN $\to$ ĐN $\to$ HCM $\to$ CT $\to$ HN có chi phí
			\[ L_2 = 40 + 50 + 30 + 110 = 230.\]
			\item Hành trình HN $\to$ HCM $\to$ ĐN $\to$ CT $\to$ HN có chi phí
			\[ L_3 = 90 + 50 + 70 + 110 = 320.\]
		\end{itemize}
		Các hành trình khác nhau nhưng đưa đến kết quả giống nhau không liệt kê ở đây. Ví dụ hành trình HN $\to$ ĐN $\to$ CT $\to$ HCM $\to$ HN cũng có kết quả giống hành trình HN $\to$ HCM $\to$ CT $\to$ ĐN $\to$ HN.\\
		So sánh các kết quả, Nam và nhóm bạn cần đổ ít nhất $230$ lít xăng để hoàn thành hành trình.
	}
\end{ex}

\begin{ex}%[Đề KSCL lần I, Sở Hưng Yên, năm học 2025-2026]%[2H2V1-4]
	Một chiếc Container được buộc vào móc $S$ của một chiếc cần cẩu bởi bốn sợi dây cáp không giãn $SA$, $SB$, $SC$, $SD$ có độ dài bằng nhau và cùng tạo với mặt phẳng $(ABCD)$ một góc bằng $30^\circ$. Chiếc cần cẩu kéo chiếc Container lên theo phương thẳng đứng. Tính độ lớn lực căng (đơn vị kN) của mỗi sợi dây cáp, biết rằng các lực căng $\overrightarrow{F}_1$, $\overrightarrow{F}_2$, $\overrightarrow{F}_3$, $\overrightarrow{F}_4$ trên mỗi sợi dây cáp có độ lớn bằng nhau và trọng lượng của chiếc Container bằng $48$ kN.
	\begin{center}
		\begin{tikzpicture}[line join=round, line cap=round, >=stealth, font=\footnotesize, scale=1]
			\clip (-2,-2.5) rectangle (3,3.5);
			%Cần cẩu
			\fill[orange!60!red](-0.85,2.24) -- (-0.17,2.5) -- (-0.17,3.14) -- (-0.84,2.87);
			\fill[orange!40!red](0.37,2.49) -- (-0.17,2.5) -- (-0.17,3.14) -- (0.34,2.95);
			\fill[orange!30!red](-0.85,2.24) -- (-0.17,2.5) -- (0.23,2.3) -- (-0.14,2.06);
			\fill[orange!70!red](-0.02,2.97) -- (0.26,3.07) -- (0.26,2.56) -- (-0.03,2.47);
			\fill[orange!40!red](0.72,2.89) -- (0.26,3.07) -- (0.26,2.56) -- (0.74,2.39);
			\fill[black!70!red](0.68,2.84) -- (0.33,2.97) -- (0.33,2.58) -- (0.69,2.45);
			\fill[orange!70!red](-0.65,2.26) -- (0.05,2.53) -- (0.05,2.64) -- (-0.64,2.37);
			\fill[orange!40!red](-0.65,2.26) -- (0.05,2.53) -- (0.83,2.26) -- (0.12,2.01);
			\fill[orange!50!red](0.05,2.65) -- (0.05,2.53) -- (0.83,2.26) -- (0.82,2.37);
			\draw[red!80!black,line width=2pt] (0.26,2.74) -- (3.36,3.88)(0.54,2.72) -- (3.41,3.75)(0.32,2.46) -- (3.76,3.65)(0.78,2.29) -- (3.83,3.29);
			\draw [red!80!black,very thick](0.82,2.32) -- (0.95,2.86) -- (1.33,2.5) -- (1.57,3.11) -- (2.1,2.75) -- (2.26,3.34) -- (2.92,3.01) -- (3.03,3.63) -- (3.64,3.25);
			\draw[red!80!black,very thick] (0.44,2.5) -- (0.72,2.9) -- (1,2.7) -- (1.25,3.11) -- (1.66,2.93) -- (1.85,3.33) -- (2.36,3.18) -- (2.59,3.6) -- (3.09,3.43) -- (3.17,3.81) -- (3.64,3.63);
			\draw [red!80!black,very thick](0.66,2.89) -- (0.82,2.81)(1.36,3.15) -- (1.52,3.07)(2.45,3.55) -- (2.62,3.46)(1.09,2.71) -- (1.48,2.53)(1.56,2.87) -- (1.99,2.69)(2.14,3.08) -- (2.59,2.87)(2.91,3.35) -- (3.28,3.13);
			%Dây cáp
			\draw[thick] (0.05,1.36)coordinate(S) -- (-1.63,-0.9)coordinate(B);
			\draw[thick]  (0.05,1.36) -- (-0.88,-0.87)coordinate(A);
			\draw[thick]  (0.05,1.36) -- (0.96,-0.74)coordinate(C);
			\draw[thick]  (0.05,1.36) -- (1.73,-0.81)coordinate(D);
			\draw[very thick]  (0.05,1.36) -- (0.06,2.26);
			%Container
			\fill[red!70!black] (-1.709,-0.879) -- (-1.64,-0.89) -- (-1.639,-0.909) -- (-1.709,-0.909)--cycle;
			\fill[red!70!white] (-1.74,-0.92) -- (0.92,-0.77) -- (0.97,-2.17) -- (-1.76,-1.94)--cycle;
			\fill[red!60!black] (0.77,-0.78) -- (0.79,-2.07) -- (0.83,-2.07) -- (0.82,-0.77);
			\fill[red!60!black] (0.6,-0.79) -- (0.61,-2.05) -- (0.66,-2.06) -- (0.65,-0.79);
			\fill[red!60!black] (0.42,-0.8) -- (0.43,-2.04) -- (0.47,-2.04) -- (0.47,-0.8);
			\fill[red!60!black] (0.25,-0.81) -- (0.25,-2.02) -- (0.3,-2.03) -- (0.3,-0.81);
			\fill[red!60!black] (0.09,-0.81) -- (0.08,-2.02) -- (0.13,-2.02) -- (0.13,-0.82);
			\fill[red!60!black] (-0.08,-0.83) -- (-0.08,-2) -- (-0.03,-2) -- (-0.03,-0.83);
			\fill[red!60!black] (-0.24,-0.84) -- (-0.25,-1.99) -- (-0.2,-1.99) -- (-0.19,-0.84);
			\fill[red!60!black] (-0.39,-0.84) -- (-0.39,-1.98) -- (-0.34,-1.98) -- (-0.34,-0.84);
			\fill[red!60!black] (-0.54,-0.85) -- (-0.55,-1.96) -- (-0.5,-1.97) -- (-0.49,-0.85);
			\fill[red!60!black] (-0.68,-0.86) -- (-0.69,-1.95) -- (-0.64,-1.96) -- (-0.63,-0.86);
			\fill[red!60!black] (-0.82,-0.87) -- (-0.83,-1.95) -- (-0.78,-1.95) -- (-0.77,-0.87);
			\fill[red!60!black] (-0.95,-0.89) -- (-0.97,-1.93) -- (-0.92,-1.94) -- (-0.91,-0.88);
			\fill[red!60!black] (-1.08,-0.89) -- (-1.1,-1.94) -- (-1.06,-1.94) -- (-1.04,-0.88);
			\fill[red!60!black] (-1.21,-0.89) -- (-1.23,-1.92) -- (-1.19,-1.92) -- (-1.17,-0.9);
			\fill[red!60!black] (-1.34,-0.91) -- (-1.36,-1.9) -- (-1.32,-1.91) -- (-1.29,-0.9);
			\fill[red!60!black] (-1.45,-0.92) -- (-1.49,-1.9) -- (-1.44,-1.9) -- (-1.42,-0.91);
			\fill[red!60!black] (-1.57,-0.92) -- (-1.6,-1.89) -- (-1.56,-1.89) -- (-1.53,-0.92);
			\fill[red!60!black] (-1.7,-0.93) -- (-1.72,-1.88) -- (-1.69,-1.89) -- (-1.66,-0.93);
			\draw[red!60!white,very thin] (0.71,-0.77) -- (0.72,-2.07)(0.54,-0.78) -- (0.55,-2.05)(0.2,-0.81) -- (0.2,-2.02)(0.36,-0.79) -- (0.37,-2.04)(0.03,-0.81) -- (0.02,-2)(-0.12,-0.83) -- (-0.13,-2)(-0.28,-0.84) -- (-0.29,-1.99)(-0.43,-0.85) -- (-0.44,-1.98)(-0.58,-0.86) -- (-0.6,-1.96)(-0.72,-0.86) -- (-0.73,-1.96)(-0.85,-0.86) -- (-0.87,-1.94)(-0.99,-0.88) -- (-1.01,-1.93)(-1.12,-0.89) -- (-1.14,-1.92)(-1.25,-0.9) -- (-1.27,-1.91)(-1.38,-0.9) -- (-1.41,-1.9)(-1.49,-0.91) -- (-1.53,-1.89)(-1.61,-0.93) -- (-1.64,-1.88);
			\fill[red!40!black] (-1.7,-0.91) -- (0.91,-0.74) -- (0.91,-0.8) -- (-1.71,-0.96)--cycle;
			\fill[red!70!white] (-1.7,-0.91) -- (0.91,-0.74) -- (0.92,-0.79) -- (-1.71,-0.95)--cycle;
			\fill[red!70!white] (-1.73,-1.87) -- (0.92,-2.07) -- (0.92,-2.17) -- (-1.73,-1.96)--cycle;
			\fill[red!70!white] (-1.78,-1.95) -- (-1.72,-1.95) -- (-1.7,-0.87) -- (-1.76,-0.87)--cycle;
			\fill[black] (-1.74,-0.88) -- (-1.72,-0.88) -- (-1.72,-0.93) -- (-1.74,-0.93)--cycle;
			\fill[black] (-1.77,-1.89) -- (-1.75,-1.89) -- (-1.75,-1.94) -- (-1.77,-1.94);
			\fill[red!60!white] (1,-0.72) -- (1.88,-0.81) -- (1.91,-2.08) -- (1.02,-2.18)--cycle;
			\draw[black!80] (1.03,-0.87) -- (1.86,-0.93)(1.04,-1.2) -- (1.86,-1.22)(1.04,-1.41) -- (1.85,-1.42)(1.04,-1.62) -- (1.87,-1.6)(1.05,-1.93) -- (1.87,-1.88);
			\draw [black!70] (1.03,-0.9) -- (1.85,-0.96)(1.03,-1.24) -- (1.86,-1.25)(1.04,-1.45) -- (1.84,-1.46)(1.05,-1.66) -- (1.86,-1.63)(1.05,-1.97) -- (1.86,-1.91);
			\draw[thick] (1,-0.78) -- (1.87,-0.85)(1.1,-2.09) -- (1.85,-2.02)(1.45,-0.81) -- (1.47,-2.05);
			\draw[black!10]  (1.24,-0.77) -- (1.27,-2.12)(1.42,-0.78) -- (1.44,-2.11) (1.52,-0.78) -- (1.54,-2.11)(1.68,-0.8) -- (1.7,-2.09)(1.26,-1.8) -- (1.12,-1.81)(1.43,-1.79) -- (1.31,-1.79)(1.53,-1.78) -- (1.64,-1.78)(1.69,-1.77) -- (1.81,-1.77)(1.17,-1.78) -- (1.17,-1.83)(1.35,-1.77) -- (1.35,-1.82)(1.61,-1.75) -- (1.61,-1.8)(1.77,-1.75) -- (1.77,-1.79);
			\fill[red!60!white] (1,-0.69) -- (1.02,-2.2) -- (1.09,-2.19) -- (1.07,-0.69)--cycle;
			\fill[red!70!white] (1,-0.69) -- (1.02,-2.2) -- (0.92,-2.19) -- (0.89,-0.7)--cycle;
			\fill[red!60!white] (1.82,-0.76) -- (1.9,-0.77) -- (1.93,-2.1) -- (1.85,-2.11)--cycle;
			\draw[fill=black] (0.93,-0.71) -- (0.96,-0.71) -- (0.96,-0.77) -- (0.93,-0.77)--cycle (0.95,-2.11) -- (0.98,-2.11) -- (0.98,-2.17) -- (0.95,-2.17)--cycle(1.05,-2.11) -- (1.08,-2.11) -- (1.08,-2.17) -- (1.05,-2.17)--cycle(1.03,-0.71) -- (1.06,-0.71) -- (1.06,-0.77) -- (1.03,-0.77)--cycle(1.85,-0.79) -- (1.88,-0.79) -- (1.88,-0.84) -- (1.85,-0.84)--cycle(1.87,-2.03) -- (1.9,-2.03) -- (1.9,-2.08) -- (1.87,-2.08)--cycle;
			\fill[black](-0.9,-1.95) -- (-0.79,-1.96) -- (-0.79,-2.01) -- (-0.9,-2);
			\fill[black](-0.31,-2) -- (-0.2,-2.01) -- (-0.2,-2.06) -- (-0.31,-2.05);
			%Véc tơ lực
			\coordinate(B1) at ($(S)!0.5!(B)$);
			\coordinate(A1) at ($(S)!0.4!(A)$);
			\coordinate(D1) at ($(S)!0.4!(D)$);
			\coordinate(C1) at ($(S)!0.55!(C)$);
			\draw[thick,->](S)--(A1);
			\draw[thick,->](S)--(B1);
			\draw[thick,->](S)--(C1);
			\draw[thick,->](S)--(D1);
			\foreach \x/\g in {A/40,B/160,C/150,D/30,S/180}
			\draw[fill=white](\x)+(\g:.3)node{\scriptsize $\x$};
			\foreach \x/\g/\i in {A1/-10/1,B1/160/2,C1/-120/3,D1/30/4}
			\draw[fill=white](\x)+(\g:.3)node{\tiny $\overrightarrow{F}_{\i}$};
		\end{tikzpicture}
	\end{center}
	\shortans{24}
	\loigiai{
		Gọi $F$ là độ lớn lực căng trên mỗi sợi dây dây cáp ($F_1=F_2=F_3=F_4=F$). \\
		Trọng lực $\vec{P}$ của Container hướng thẳng xuống dưới, có độ lớn $P = 48$ kN.
		\begin{center}
			\begin{tikzpicture}[scale=1, font=\footnotesize,line join=round, line cap=round, >=stealth]
				\def\h{3}
				\path (0,0) coordinate (A)
				(-1.5,-1.2) coordinate (B)
				++(4,0) coordinate (C)
				($(A)+(C)-(B)$) coordinate(D)
				($(A)!1/2!(C)$) coordinate(O)
				++(0,\h) coordinate(S);
				\foreach \i/\j/\k in {A/1/-10, B/2/160, C/3/10, D/4/20} 
				{\draw[->, red] (S) -- ($(S)!0.6!(\i)$) coordinate(F\j) node[shift={(\k:4mm)}]{$\overrightarrow{F}_\j$};}
				\draw (S)--(B)--(C)--(D)--(S)--(C);
				\draw[dashed] (S)--(A)--(B) (A)--(D) (S)--(O) (A)--(C) (B)--(D);
				\foreach \diem/\pos in {A/160, B/-135, C/-45, D/0, S/90, O/-90} \fill (\diem) node[shift={(\pos:0.3)}] {$\diem$} circle(1pt);
				\path (0,0) node{\hypersetup{hidelinks}\href{tpSzVVc}{ }};
			\end{tikzpicture}
		\end{center}
		Khi đó ta có các lực $\overrightarrow{F}_1$, $\overrightarrow{F}_2$, $\overrightarrow{F}_3$, $\overrightarrow{F}_4$ lần lượt cùng hướng với các véc-tơ $\overrightarrow{SA}$, $\overrightarrow{SB}$, $\overrightarrow{SC}$, $\overrightarrow{SD}$.\\
		Do đó tồn tại $k>0$ để
		$\overrightarrow{F}_1=k\overrightarrow{SA}$, $\overrightarrow{F}_2=k\overrightarrow{SB}$, $\overrightarrow{F}_3=k\overrightarrow{SC}$, $\overrightarrow{F}_4=k\overrightarrow{SD}$.\\
		Gọi $\overrightarrow{P}$ là véc-tơ trọng lực tác động lên Container.\\
		Ta có $$\overrightarrow{P} = \overrightarrow{F}_1 + \overrightarrow{F}_2 + \overrightarrow{F}_3 + \overrightarrow{F}_4 = k\left(\overrightarrow{SA} + \overrightarrow{SB} + \overrightarrow{SC} + \overrightarrow{SD}\right) = 4k\overrightarrow{SO}.$$
		Suy ra $F=\left|\overrightarrow{F}_1\right| = kSA$, $P=\left|\overrightarrow{F}_1 + \overrightarrow{F}_2 + \overrightarrow{F}_3 + \overrightarrow{F}_4\right| = 4kSO$.
		Tam giác $SAO$ vuông tại $O$, góc $\widehat{SAO}$ là góc giữa $SA$ với $(ABCD)$, do đó $SO=SA\cdot \sin30^\circ$ nên
		\[P=4kSO = 4k\cdot SA\cdot \sin30^\circ = 4\cdot F\cdot \sin30^\circ \Rightarrow F=\dfrac{P}{4\sin30^\circ} = \dfrac{48}{2}=24 \text{ (kN)}.\]
		Vậy độ lớn lực căng của mỗi sợi dây cáp là $24$ kN.
	}
\end{ex}
\Closesolutionfile{ans}
% \indapan{4}{ans/ans-2-KS-SoHungYen-26}


